\section{Teekay Tankers (NYSE: TNK)}
\label{sec:tnk}

Kurs \je{\TnkKurs}{USD} am \TnkKursDatum{}\QW{Yahoo Finance, Kursreihe TNK}{quelle:kurse}{19.09.2026},
\Mio{\TnkAktienA}{} Class-A-Aktien\Q{TnkZF}{2} und \Mio{\TnkAktienB}{} Class-B-Aktien\Q{TnkZF}{66},
B\"orsenwert \Mio{\TnkBoersenwert}{USD}.

\subsection{Eigent\"umer und Directors' Dealings}

Teekay Corporation h\"alt \Pz{\TnkTeekayAnteil} aller Class-A- und Class-B-Aktien und
s\"amtliche Class-B-Aktien; diese tragen f\"unf Stimmen je Aktie, begrenzt auf einen
Stimmenanteil von 49\,\%\Q{TnkZF}{66}. Die Gesellschaft ist damit von ihrem
Hauptaktion\"ar kontrolliert, ohne dass dieser die Kapitalmehrheit h\"alt.

Directors' Dealings: Teekay Tankers ist ein ausl\"andischer Emittent, der nur an der NYSE
notiert. F\"ur ihn gilt weder die Meldepflicht des SEC-Formulars~4 noch Art.~19 der
EU-Markt\-miss\-brauchs\-ver\-ord\-nung. Es gibt keine \"offentliche Quelle; das ist eine
strukturelle L\"ucke, keine unvollendete Suche.

\subsection{Wettbewerbsposition}

Teekay Tankers beschreibt den eigenen Markt als \enquote{the highly competitive
international tanker market}\Q{TnkZF}{7}. Einen Vorteil, der sich gegen\"uber anderen
Reedereien in einer Zahl zeigen w\"urde, behauptet der Bericht nicht.

\bk{Das ist ehrlich und f\"ur die Bewertung entscheidend: Ohne Vorteil folgt der Zufluss
der Frachtrate, und die Frachtrate folgt dem Verh\"altnis von Tonnage und Nachfrage. Die
Rechnung muss deshalb den ganzen Zyklus sehen, nicht sein bestes Jahr.}

\subsection{Mehrjahresreihe}

\reihentabelle{tnk}{USD}{tnk}

Median \TnkFcfJahre{}: \Mio{\TnkFcfMedian}{USD}; 2025: \Mio{\TnkFcfVoll}{USD}, \"uber
\Pz{\TnkFcfRang} der Jahre. Der Zufluss 2025 enth\"alt die Schiffsk\"aufe des Jahres;
Verkaufserl\"ose f\"ur Schiffe sind nach der Methode nicht angerechnet. F\"ur die Jahre
2017 bis 2019 stammen die Werte aus der \"Ubersichtsseite des 20-F 2019, weil dessen
Abschlussseiten keine Seitenzahlen tragen.

\annahme{Die Class-B-Aktienzahl ist die zum 1.~M\"arz 2026 genannte, die Class-A-Zahl die
zum 31.~Dezember 2025; die beiden Stichtage liegen zwei Monate auseinander.}

\subsection{Kursverlauf}

\kgvtabelle{tnk}{USD}

\bk{Das KGV liegt mit \TnkKgvHeute{} \"uber \Pz{\TnkKgvRang} der erfassten Jahre (Median
\TnkKgvMedian). Seit dem Jahresende 2025 hat sich der Kurs fast verdoppelt, ohne dass ein
neues Jahresergebnis vorliegt: Der Markt preist h\"ohere Frachtraten ein, die in den
Zahlen noch nicht stehen.}

\subsection{Votum}

\schwellentabelle{tnk}{USD}{tnk}
\umkehrtabellen{tnk}

\vd{\textbf{Nicht kaufen} zu \je{\TnkKurs}{USD}. Schwelle auf zehn Jahre
\je{\TnkSchwelleMEDZehn}{USD} (Median), \je{\TnkSchwelleVOLLZehn}{USD} (2025); bei
\Pz{\HuerdeLang} \je{\TnkSchwelleMEDZehnLang}{USD}. Der Kurs ist das
\TnkVielfachesMED-fache. Halter: der Kurs liegt \"uber dem Buchwert von
\je{\TnkBuchwertJeAktie}{USD} je Aktie; Gewinne zu sichern ist vertretbar. Nichthalter:
nicht einsteigen.}

\vd{\emph{Begr\"undung.} Der Preis verlangt \TnkWachstumMEDZehn{} Wachstum p.\,a.\ \"uber
zehn Jahre. Eine realisierte Rate l\"asst sich nicht angeben, weil die ersten Jahre der
Reihe im Mittel negativ sind; die Reihe zeigt Zyklen, keinen Trend.}

\vd{\emph{Was das Votum \"andert.} Ein Kurs in der N\"ahe des Buchwerts je Aktie; eine
Ver\"offentlichung, die die H\"ohe der Frachtraten 2026 im Zufluss zeigt.}

\vd{\emph{Was dagegen spricht.} Im oberen Teil des Zyklus verdient eine Reederei in
einem Jahr, was sie in schwachen Jahren nicht verdient; wer die Frachtrate richtig
einsch\"atzt, verdient an genau dieser Spanne. Diese Rechnung sch\"atzt die Frachtrate
nicht und kann den Zeitpunkt nicht beurteilen.}
